\documentclass[10pt,letterpaper]{article}
\usepackage[spanish,es-nodecimaldot]{babel}
\usepackage[T1]{fontenc}
\usepackage[utf8]{inputenc}
\usepackage[margin=2.0cm]{geometry}
\usepackage{array,tabularx,booktabs}
\usepackage[table]{xcolor}
\usepackage{enumitem}
\usepackage{hyperref}
\usepackage{fancyhdr}
\usepackage{microtype}
\hypersetup{colorlinks=true,linkcolor=teal!60!black,urlcolor=blue!60!black,citecolor=teal!60!black}
\definecolor{verde}{HTML}{174A6E}
\definecolor{azul}{HTML}{174A6E}
\definecolor{palido}{HTML}{EAF2F7}
\newcommand{\separa}{\hspace{0.25em}\raisebox{0.25ex}{\rule{0.38ex}{0.38ex}}\hspace{0.25em}}
\renewcommand{\arraystretch}{1.16}
\setlength{\parindent}{0pt}
\setlength{\parskip}{5pt}
\pagestyle{fancy}
\fancyhf{}
\fancyhead[L]{UTEQ \separa Ingeniería de Requisitos (ISR-401)}
\fancyhead[R]{GA-TL-02 \separa FORM}
\fancyfoot[L]{ISR-401 \separa Ingeniería de Requisitos}
\fancyfoot[R]{Página \thepage}
\renewcommand{\headrulewidth}{0.4pt}
\renewcommand{\footrulewidth}{0pt}
\newcommand{\titulo}[1]{\section*{\textcolor{verde}{#1}}}
\newcommand{\subtitulo}[1]{\subsection*{\textcolor{azul}{#1}}}
\begin{document}
\begin{center}
{\Large\bfseries\textcolor{verde}{Guía de entrevista semiestructurada}}\\[3pt]
{\large GA-TL-02 \separa FORM \separa Taller de elicitación y Design Thinking}
\end{center}
\textbf{Integrantes:} Alvarez Rubio Rai Eduardo; Briones Solorzano William Adrian; Rivera Rodriguez Joseph Jeremy\\
\textbf{Responsable de esta parte:} Alvarez Rubio Rai Eduardo\\
\textbf{Asignatura:} ISR-401 \separa Ingeniería de Requisitos \separa UTEQ\\
\textbf{Proyecto:} Sistema de gestión para un gimnasio\\
\textbf{Persona entrevistada:} Propietario o administrador del gimnasio.

\titulo{1. Objetivo de la entrevista}
Conocer cómo se organiza el gimnasio, qué actividades realizan sus responsables y usuarios, qué herramientas utilizan y qué situaciones les resultan difíciles en el trabajo diario. También se busca identificar expectativas y condiciones que deberían considerarse antes de proponer un sistema de gestión. La intención es comprender primero el contexto, no llegar con una solución decidida de antemano.

\titulo{2. Preparación de la entrevista}
\begin{tabularx}{\linewidth}{@{}p{4.2cm}X@{}}
\rowcolor{palido}\textbf{Duración prevista} & 20 a 30 minutos para el levantamiento real. El roleplay de clase se ajustará a 10 minutos.\\
\textbf{Modalidad} & Presencial o virtual, según la disponibilidad de la persona entrevistada.\\
\textbf{Entrevistado sugerido} & Propietario o administrador del gimnasio.\\
\textbf{Materiales} & Esta guía, protocolo de consentimiento, notas y, solo con autorización expresa, dispositivo de grabación.\\
\textbf{Entrevistador} & Alvarez Rubio Rai Eduardo.\\
\textbf{Registro} & Un integrante tomará notas de ideas, ejemplos, dudas y términos utilizados por la persona entrevistada.\\
\end{tabularx}

\titulo{3. Apertura y encuadre}
\textit{Presentación sugerida:}

Muchas gracias por recibirnos. Somos estudiantes de Ingeniería de Software y estamos explorando la posibilidad de desarrollar un sistema relacionado con la gestión de un gimnasio. Antes de pensar en pantallas o funcionalidades, nos gustaría conocer cómo trabajan actualmente, qué les funciona bien y qué cosas les gustaría mejorar. No hay respuestas correctas o incorrectas; nos interesa conocer su experiencia tal como es. La información se utilizará únicamente para el trabajo académico y seguiremos lo indicado en el consentimiento informado.

\titulo{4. Preguntas organizadas por bloques}

\subtitulo{Bloque A. Conocer el gimnasio y su funcionamiento}
\begin{enumerate}[leftmargin=*,itemsep=5pt]
\item Para comenzar, ¿podría contarnos un poco sobre el gimnasio y las actividades o servicios que ofrecen?
\item ¿Cómo se distribuyen las responsabilidades entre las personas que trabajan aquí?
\item ¿Cómo transcurre un día normal en el gimnasio, desde que abre hasta que termina la jornada?
\item ¿Qué tareas administrativas o de atención realizan con mayor frecuencia?
\end{enumerate}

\subtitulo{Bloque B. Inscripción y relación con los socios}
\begin{enumerate}[leftmargin=*,itemsep=5pt,resume]
\item Cuando una persona quiere inscribirse por primera vez, ¿qué pasos siguen para atenderla y registrarla?
\item ¿Qué información suelen solicitar a los nuevos socios y para qué la utilizan?
\item ¿Cómo registran los planes, las membresías, sus fechas de inicio y vencimiento?
\item Cuando un socio desea cambiar de plan, pausar su membresía o retirarse, ¿cómo gestionan esa solicitud?
\end{enumerate}

\subtitulo{Bloque C. Pagos y control administrativo}
\begin{enumerate}[leftmargin=*,itemsep=5pt,resume]
\item ¿Cómo llevan actualmente el control de los pagos y de las mensualidades pendientes?
\item ¿Qué hacen para saber qué membresías están vigentes y cuáles están próximas a vencer?
\item ¿Cómo registran los comprobantes o verifican que un pago se haya realizado correctamente?
\item ¿Preparan reportes sobre ingresos, socios activos u otros datos? ¿Cómo los elaboran?
\end{enumerate}

\subtitulo{Bloque D. Asistencia, servicios y comunicación}
\begin{enumerate}[leftmargin=*,itemsep=5pt,resume]
\item ¿Cómo controlan la asistencia de los socios cuando ingresan al gimnasio?
\item ¿Manejan horarios de clases, turnos, entrenadores o reservas? ¿Cómo se organizan?
\item ¿Por qué medios se comunican con los socios para informar cambios, promociones o novedades?
\item ¿Qué consultas o solicitudes reciben con más frecuencia por parte de los socios?
\end{enumerate}

\subtitulo{Bloque E. Dificultades y expectativas}
\begin{enumerate}[leftmargin=*,itemsep=5pt,resume]
\item De todas las tareas que hemos conversado, ¿cuál suele tomar más tiempo o causar más inconvenientes? ¿Podría compartir un ejemplo?
\item ¿En qué momentos se producen confusiones, errores o reclamos, si es que ocurren?
\item ¿Qué información les cuesta encontrar cuando necesitan tomar una decisión o atender a un socio?
\item Si pudieran mejorar una sola parte de la gestión del gimnasio, ¿cuál elegirían y por qué?
\item Pensando en una herramienta informática, ¿qué tarea les gustaría que facilitara primero y qué resultado esperarían obtener?
\item ¿Hay alguna condición que debamos tener en cuenta para que una herramienta así sea práctica para ustedes, por ejemplo, el equipo disponible, la conexión a internet o el tiempo para aprender a usarla?
\end{enumerate}

\titulo{5. Preguntas de profundización}
No es necesario formular todas estas preguntas. El entrevistador puede utilizarlas cuando una respuesta sea muy general o cuando sea útil conocer un caso concreto.

\begin{itemize}[leftmargin=*,itemsep=3pt]
\item ¿Podría explicarme un poco más cómo sucede eso?
\item ¿Recuerda la última vez que ocurrió? ¿Qué hicieron para resolverlo?
\item ¿Quién suele encargarse de esa tarea?
\item ¿Qué información necesita tener a mano para completar ese proceso?
\item ¿Qué pasa después de ese paso?
\item ¿Qué sería diferente para usted si ese proceso funcionara mejor?
\end{itemize}

\titulo{6. Cierre de la conversación}
\textit{Cierre sugerido:}

Muchas gracias por compartir cómo funciona el gimnasio y por contarnos su experiencia. Lo que nos ha explicado nos ayudará a entender mejor las necesidades antes de definir qué debería hacer el sistema. Antes de terminar, ¿hay algo importante que no le hayamos preguntado y que considere necesario que conozcamos? Si al revisar nuestras notas nos surge alguna duda, ¿sería posible consultarle nuevamente por el medio que usted prefiera?

\titulo{7. Recomendaciones para quien entrevista}
\begin{itemize}[leftmargin=*,itemsep=3pt]
\item Escuchar sin interrumpir y permitir que la persona explique con sus propias palabras.
\item Evitar sugerir respuestas o presentar como hechos problemas que todavía no se han confirmado.
\item Pedir ejemplos concretos cuando se mencionen dificultades, excepciones o tareas repetitivas.
\item Mantener el orden de los bloques como orientación; se puede cambiar el orden si la conversación lo requiere.
\item Diferenciar lo que la persona dijo de las interpretaciones o ideas del equipo.
\item Anotar términos desconocidos para aclararlos, en lugar de asumir su significado.
\end{itemize}

\titulo{8. Registro de la entrevista}

\begin{tabularx}{\linewidth}{@{}p{4.0cm}X@{}}
\textbf{Fecha y modalidad} & \rule{0pt}{0.8cm}\\
\textbf{Persona y cargo} & \rule{0pt}{0.8cm}\\
\textbf{Procesos descritos} & \rule{0pt}{1.2cm}\\
\textbf{Dificultades mencionadas} & \rule{0pt}{1.2cm}\\
\textbf{Necesidades o expectativas} & \rule{0pt}{1.2cm}\\
\textbf{Dudas por confirmar} & \rule{0pt}{1.2cm}\\
\end{tabularx}

\vfill
\begin{center}
\small\textcolor{verde}{Documento de trabajo académico \separa GA-TL-02 \separa FORM}
\end{center}
\end{document}
